While deep neural networks (DNNs) have demonstrated remarkable capabilities, achieving human-like or even superior performance across a wide range of tasks, their robustness is often compromised by their susceptibility to input perturbations \cite{szegedy}. This vulnerability has catalyzed the formal verification community to develop various methodologies, each presenting a unique trade-off between completeness and computational efficiency~\cite{Marabou,Reluplex,deeppoly}. 
This surge in innovation has also led to the establishment of competitions such as VNNComp~\cite{VNNcomp}, which systematically evaluate the performance of neural network verification tools. Notable examples include NNenum~\cite{nnenum}, Marabou~\cite{Marabou,Marabou2}, and PyRAT~\cite{pyrat}, as well as frameworks such as MnBAB~\cite{ferrari2022complete} (which builds upon ERAN~\cite{deeppoly} and PRIMA~\cite{prima}) and $\alpha,\beta$-CROWN~\cite{crown,xu2020fast}, based on branch-and-bound methodology \cite{BaB}.

However, these verification tools focus on {\em local} robustness: given a DNN, an image, and a small neighborhood around this image, they verify whether all images in the neighborhood are assigned the same classification. This neighborhood is provided by a maximal perturbation of the input image, often an $L_\infty$-norm constraint. Under this metric, every subpixel of the input image can vary in a very small range, typically $\frac{2}{255}$ (that is, 2 levels of gray/blue/red/green). 
Although it is not necessarily the most meaningful perturbation, $L_\infty$ is the usual choice because it is perfectly linear. 
%(subpixel perturbations are independent), which is easier to verify.

%- to line break
Crucially, however, these verification tools for local robustness are too computationally intensive for real-time decision-making pipelines. 
Consider an autonomous vehicle processing a live video feed: images of the feed cannot be certified as robust in a few ms on embedded hardware. This prevents the development of safety filters that could skip non-robust images and only consider certified images.

\smallskip

In this paper, we achieve {\em real-time} certification 
using {\em global} robustness bounds \eca{that remain valid across the entire input space}. 
%(that do not restrict the verification process to the neighborhood of a fixed input)
Our approach follows a two-step procedure:

\textit{Global Bound Analysis (Offline).} We compute global bounds on the maximum possible \eca{difference} between output values of different decision classes due to a perturbation, following an approach similar to VHAGaR~\cite{vhagar}. 
%
Specifically, for any two decision classes $C$ and $D$ and perturbation $\varepsilon$, we compute an 
upper bound 
$\bar{\beta}^\varepsilon_{C,D}$ 
on $\max_{I,I', |I-I'| \leq \epsilon}(
x_C - x'_C + x'_D - x_D)$, where $x_{C,D},x'_{C,D}$ are the output values of class $C,D$ from image $I$ and $I'$. 
%
Unlike local robustness methods, which require $k$ separate, computationally-expensive calls to certify $k$ images, the global upper bound $\bar{\beta}^\varepsilon_{C,D}$ is computed offline once per network and remains valid across the entire input space. 

% NOTE: change eq. to align with what is used later on
\textit{Real-Time Verification (Online).} \eca{During deployment, after the DNN's standard inference phase for a given input image $I$}, we identify the class $C$ with the highest output value $x_C$ and check whether for every other class \eca{$D \neq C$, 
$x_D +\bar{\beta}^{\varepsilon}_{C,D} \le x_C$}. Thereby, we can \eca{formally} certify that the image $I$ is robust for \eca{any} perturbation $\varepsilon$. This check typically requires only a few dozen CPU instructions (under 1 ms on a single embedded CPU core). When a frame fails this verification, one could either skip it or revert to a safer, degraded mode until a trustworthy, robust frame is received.

Our main contributions address the challenge of computing {\em global bounds} $\bar{\beta}^\varepsilon_{C,D}$ for classes $(C,D)$:

\begin{enumerate}
	\item \emph{AutoEncoders for ID Projection.} \eca{Global bounds for vanilla DNNs} are particularly pessimistic because inputs far away from the training dataset, i.e., not {\em In Distribution (ID)}, need to be accounted for. \eca{Since DNNs typically behave erratically on such inputs \cite{OOD}, this results in large lower bounds on ${\beta}^\varepsilon_{C,D}$ (see Section \ref{sec.exp}), thus {\em no} technique can produce manageable upper bound $\bar{\beta}^\varepsilon_{C,D}$.} While \eca{explicitly} {\em characterizing} ID inputs to restrict \eca{the analysis} to them is impractical, we propose an efficient alternative: learn from ID samples an autoencoder (AE), and then use the AE-DNN architecture (Fig.~\ref{fig:AE-DNN} in Section 3) to first project the input on a small latent space \eca{whose size is selected to ensure that AE-DNN accuracy is within 1\% of the original DNN accuracy.} We experimented with Linear AE (LAE)\eca{; while the encoding may not be as efficient as with non-Linear AE, they significantly simplify bound computation.} This results in latent spaces $>30\times$ smaller than full spaces across 3 datasets, leading to LAE-DNNs' bounds 3--15$\times$ tighter than those of standard DNNs.
	%and LAE-DNNs bounds $\bar{\beta}^\varepsilon_{C,D}$ computed are 3--15$\times$ smaller than for DNNs.

	\item \emph{Diff MILP Encoding.} \eca{We introduce a novel {\em Diff} MILP encoding to enable efficient bound computation on the global robustness problem.} In this model, the variables are the values of the perturbed neurons, as well as the differences between the original and the perturbed neuron values (called the {\em diff variables}, introduced in \cite{diff}). We study and encode how the {\em diff variables} evolve after passing through a ReLU (Prop.~\ref{Prop2}), see Section \ref{s.diff}. By comparison, the {\em classical} MILP model \cite{MILP} employed in \cite{vhagar,lipshitz,ITNE} computes the {\em diff variables} as the difference between 2 larger values. Both encodings are {\em asymptotically} exact, but the {\em Diff} model is more accurate in practice. Further, from the {\em Diff MILP model}, which is asymptotically exact, we develop two abstract MILP models, which are more efficient but also asymptotically less accurate than {\em Diff}, namely  the {\em 2b+1} and  the {\em 1b} models.
	


\iffalse
	\item We develop a novel {\em Diff MILP encoding} for the global robustness problem, 
	where the variables are the values of the perturbed neurons, as well as the difference between the original and the perturbed neuron values (called the {\em diff variables}, introduced in \cite{diff}). We study and encode how the {\em diff variables} evolve after passing through a ReLU (Prop.~\ref{Prop2}), see Section \ref{s.diff}. 
	Compared with the {\em classical MILP model} \cite{MILP} employed in 
	\cite{vhagar,lipshitz,ITNE}, which considers the input and the perturbation instead of the {\em diff variable}, we keep the same number (2) of binary variables per ReLU.
	%the results are more accurate for the same runtime.
	%the linear relaxation is much more accurate, as each {\em diff variable} can be bounded after a ReLU as a function of the value of the {\em diff variables} before the ReLU, whereas the linear relaxations of the classical model is extremely inaccurate, as the variables are independent of each other. This was observed in \cite{lipshitz,ITNE}, and constraints encoding
	%a part of the linear relaxation of our "2v" model were added explicitly. 
	%Experimentally, 
	However, our {\em Diff MILP model} is more efficient, one reason being 
	the accuracy of its linear relaxation.
	%The number of variables a priori doubles compared with local robustness, from each neuron value in the perturbed image to each neuron value in the perturbed image {\em and} in the original image, as the original image is no more fixed. 
	%Recall that the worst case complexity of MILP is exponential in the number of binary variables \cite{DivideAndSlide}. 
	%A straightforward MILP model would be to use the  for each of these variables, as in \cite{vhagar,lipshitz}. The main issue with the classical model is that its linear relaxations is extremely inaccurate, as the variables are independent of each other. Instead, we develop another exact MILP model, 
\fi

	%Namely, the {\em 2b+1} model, accurate on the {\em diff variables} but abstract on the perturbed variables; while the {\em 1b} model has a unique binary variable per ReLU, only considering the {\em diff variables}.

\end{enumerate}


In terms of perturbations, we consider conjunctions of $L_\infty$- and $L_1$-norms, which allow to accurately describe perturbations. For instance, "each subpixel is perturbed by at most $\frac{50}{255}$ ($L_\infty$) and the sum of the absolute value of perturbations over all subpixels is at most $1$" ($L_1$-perturbation). While $L_1$ perturbations are not linear (because of the absolute values), and thus they are seldom used, we show in Section~\ref{s.L1} how to use them as a perturbation in the MILP model without incurring any expensive binary variables (only cheap linear variables are necessary). 

Experimentally, the {\em Diff MILP model} computes upper bounds
$\bar{\beta}^{\varepsilon}_{C,D}$ which reduces the gap to the lower bound compared with optimized versions of the classical MILP encoding (see Table \ref{table.classical}). This \eca{includes improvement over method that reduce} the number of binary variables (VHAGaR \cite{vhagar}) or incorporate linear constraints from the {\em diff variables} (ITNE \cite{ITNE}).
The abstractions {\em 2b+1} and {\em 1b} variants offer different trade-offs, 
reaching better bounds than the {\em Diff MILP} model when the instance is very complex, or when runtime is limited (Table \ref{table.L1}).  
%Further, using PCA reduces the upper bound $\bar{\beta}^{\varepsilon}_{C,D}$ by $3$ to $20$ times.
Overall, using both LAE-DNNs and {\em Diff} MILP and abstractions models, enables the real-time certification of $\geq 70\%$ of fresh images for an $L_1$-perturbation of $1,4,2$ over the MNIST, Fashion-MNIST, and CIFAR-10 datasets, respectively, see Table \ref{table.cert}. The online process only adds \eca{about} 0.5 ms of latency per image, and \eca{approximately} 2000 images/second can be treated per CPU core.

This is a Concept and Feasibility paper. 
Our techniques are tested on several DNNs and datasets, but DNN size is limited (though concrete examples are tackled). 
Linear AE and Linear and ReLU activation functions are the ones we experimented on. We strongly believe that the methodology is novel, impactful, and not intrinsically limited to these restrictions, that will eventually be lifted.